\documentclass[aspectratio=169,11pt]{beamer}

\usepackage[T1]{fontenc}
\usepackage[utf8]{inputenc}
\usepackage[french]{babel}
\usepackage{amsmath,amssymb}
\usepackage{lmodern}

\definecolor{titleblue}{RGB}{50,55,185}
\definecolor{footerleft}{RGB}{70,65,190}
\definecolor{footercenter}{RGB}{135,125,220}
\definecolor{footerright}{RGB}{185,180,235}

\setbeamercolor{normal text}{fg=black,bg=white}
\setbeamercolor{frametitle}{fg=titleblue,bg=white}
\setbeamerfont{frametitle}{size=\Large,series=\normalfont}
\setbeamertemplate{navigation symbols}{}

\setbeamertemplate{itemize item}{\color{titleblue}\small$\bullet$}
\setbeamertemplate{itemize subitem}{\color{titleblue}\scriptsize$\bullet$}

\setbeamertemplate{footline}{
    \leavevmode
    \hbox{
        \begin{beamercolorbox}[wd=.33\paperwidth,ht=2.8ex,dp=1.2ex,center]{author in head/foot}
            \color{white}\scriptsize Monte Carlo Simulation
        \end{beamercolorbox}
        \begin{beamercolorbox}[wd=.34\paperwidth,ht=2.8ex,dp=1.2ex,center]{title in head/foot}
            \color{white}\scriptsize Pricing de produits dérivés
        \end{beamercolorbox}
        \begin{beamercolorbox}[wd=.33\paperwidth,ht=2.8ex,dp=1.2ex,center]{date in head/foot}
            \color{white}\scriptsize \insertframenumber{} / \inserttotalframenumber
        \end{beamercolorbox}
    }
}

\setbeamercolor{author in head/foot}{bg=footerleft,fg=white}
\setbeamercolor{title in head/foot}{bg=footercenter,fg=white}
\setbeamercolor{date in head/foot}{bg=footerright,fg=white}

\begin{document}

%------------------------------------------------
\begin{frame}
    \centering
    \vfill

    {\color{titleblue}\Huge
        \textbf{Modélisation Monte Carlo en C++}
    }

    \vfill
\end{frame}

%------------------------------------------------
\begin{frame}
    \frametitle{Sommaire}

    \begin{itemize}
        \item Cadre mathématique de la méthode Monte Carlo

        \vspace{0.3cm}

        \item Simulation du sous-jacent sous différents modèles

        \vspace{0.3cm}

        \item Pricing d'un call européen pour chaque modèle

        \vspace{0.3cm}

        \item Pricing de produits exotiques

        \vspace{0.3cm}

        \item Méthodes de réduction de variance

        \vspace{0.3cm}

        \item Parallélisation avec OpenMP
    \end{itemize}

\end{frame}

%------------------------------------------------
\begin{frame}
    \frametitle{Cadre de valorisation}

    On se place dans un marché financier vérifiant :

    \vspace{0.4cm}

    \begin{itemize}
        \item \textbf{Absence d'arbitrage} :
        il n'existe pas de stratégie autofinancée permettant un gain certain sans risque.

        \vspace{0.25cm}

        \item \textbf{Marché complet} :
        tout payoff atteignable peut être répliqué par une stratégie de portefeuille.

        \vspace{0.25cm}

        \item \textbf{Conséquence} :
        il existe une unique mesure risque-neutre $\mathbb{Q}$.
    \end{itemize}

    \vspace{0.4cm}

    Le prix d'un produit dérivé est donné par l'espérance actualisée
    de son payoff sous $\mathbb{Q}$.

\end{frame}

%------------------------------------------------
\begin{frame}
    \frametitle{Valorisation risque-neutre}

    Soit $(V_t)_{t\geq 0}$ le prix d'un produit dérivé
    écrit sur un sous-jacent $(S_t)_{t\geq 0}$.

    \vspace{0.3cm}

    Son payoff peut dépendre de toute la trajectoire du sous-jacent :

    \[
        \Phi\left((S_u)_{t\leq u\leq T}\right).
    \]

    Sous l'unique mesure risque-neutre $\mathbb{Q}$ :

    \[
        \boxed{
        V_t
        =
        \mathbb{E}^{\mathbb{Q}}
        \left[
            \frac{B_t}{B_T}
            \Phi\left((S_u)_{t\leq u\leq T}\right)
            \,\middle|\,
            \mathcal{F}_t
        \right].
        }
    \]

\end{frame}

%------------------------------------------------
\begin{frame}
    \frametitle{Actualisation du payoff}

    L'actif sans risque $(B_t)_{t\geq 0}$ est défini par :

    \[
        \boxed{
        B_t
        =
        \exp\left(
            \int_0^t r_s\,ds
        \right).
        }
    \]

    Si le taux d'intérêt $r$ est constant :

    \[
        B_t=e^{rt}
        \qquad\text{et}\qquad
        \frac{B_t}{B_T}=e^{-r(T-t)}.
    \]

    La formule devient :

    \[
        \boxed{
        V_t
        =
        e^{-r(T-t)}
        \mathbb{E}^{\mathbb{Q}}
        \left[
            \Phi\left((S_u)_{t\leq u\leq T}\right)
            \,\middle|\,
            \mathcal{F}_t
        \right].
        }
    \]

\end{frame}

%------------------------------------------------
\begin{frame}
    \frametitle{Modélisation Monte Carlo}

    On estime le prix $V_t$ en réalisant $N$ simulations indépendantes
    de la trajectoire du sous-jacent :

    \[
        (S_u^{(1)})_{t\leq u\leq T},
        \ldots,
        (S_u^{(N)})_{t\leq u\leq T}.
    \]

    \vspace{0.4cm}

    L'estimateur Monte Carlo du prix est :

    \[
        \boxed{
        \widehat{V}_N
        =
        e^{-r(T-t)}
        \frac{1}{N}
        \sum_{i=1}^{N}
        \Phi\left((S_u^{(i)})_{t\leq u\leq T}\right).
        }
    \]

\end{frame}

%------------------------------------------------
\begin{frame}
    \frametitle{Erreur standard de l'estimateur}

    On note les payoffs simulés :

    \[
        Y_i
        =
        \Phi\left((S_u^{(i)})_{t\leq u\leq T}\right).
    \]

    La moyenne empirique est :

    \[
        \overline{Y}_N
        =
        \frac{1}{N}
        \sum_{i=1}^{N}
        Y_i.
    \]

    La variance empirique est estimée par :

    \[
        \widehat{\sigma}^{\,2}
        =
        \frac{1}{N-1}
        \sum_{i=1}^{N}
        \left(Y_i-\overline{Y}_N\right)^2.
    \]

    \[
        \boxed{
        \widehat{\mathrm{SE}}
        =
        e^{-r(T-t)}
        \frac{\widehat{\sigma}}{\sqrt{N}}.
        }
    \]

\end{frame}

%------------------------------------------------
\begin{frame}
    \frametitle{Modèle de Black--Scholes et simulation}

    Sous la mesure risque-neutre $\mathbb{Q}$, le sous-jacent vérifie :

    \[
        \boxed{
        dS_t
        =
        rS_t\,dt
        +
        \sigma S_t\,dW_t^{\mathbb{Q}}.
        }
    \]

    Entre les dates $t$ et $T$ :

    \[
        S_T
        =
        S_t
        \exp\left[
            \left(r-\frac{\sigma^2}{2}\right)(T-t)
            +
            \sigma\sqrt{T-t}\,Z
        \right],
        \qquad
        Z\sim\mathcal{N}(0,1).
    \]

\end{frame}

%------------------------------------------------
\begin{frame}
    \frametitle{Call européen}

    Pour un call européen, le payoff dépend seulement de $S_T$ :

    \[
        \boxed{
        \Phi(S_T)
        =
        (S_T-K)^+
        =
        \max(S_T-K,0).
        }
    \]

    Son prix à l'instant $t$ est :

    \[
        \boxed{
        V_t
        =
        e^{-r(T-t)}
        \mathbb{E}^{\mathbb{Q}}
        \left[
            (S_T-K)^+
            \,\middle|\,
            \mathcal{F}_t
        \right].
        }
    \]

\end{frame}

%------------------------------------------------
%------------------------------------------------
\begin{frame}
    \frametitle{Algorithme de Welford : mise à jour}

    L'algorithme de Welford calcule la moyenne et la variance
    empiriques en une seule passe, sans stocker tous les payoffs.

    \vspace{0.3cm}

    Initialisation :

    \[
        n=0,
        \qquad
        \overline{Y}_0=0,
        \qquad
        M_{2,0}=0.
    \]

    Pour chaque nouveau payoff $Y_n$, avec $n\geq 1$ :

    \[
        \delta_n
        =
        Y_n-\overline{Y}_{n-1},
    \]

    \[
        \overline{Y}_n
        =
        \overline{Y}_{n-1}
        +
        \frac{\delta_n}{n},
    \]

    \[
        M_{2,n}
        =
        M_{2,n-1}
        +
        \delta_n
        \left(
            Y_n-\overline{Y}_n
        \right).
    \]

\end{frame}

%------------------------------------------------
%------------------------------------------------
\begin{frame}
    \frametitle{Algorithme de Welford : variance et erreur standard}

    Après avoir traité les $N$ payoffs simulés, l'algorithme fournit :

    \[
        \overline{Y}_N
        =
        \frac{1}{N}
        \sum_{i=1}^{N}Y_i,
        \qquad
        \widehat{\sigma}_N^2
        =
        \frac{M_{2,N}}{N-1}.
    \]

    \vspace{0.3cm}

    L'estimateur Monte Carlo du prix est alors :

    \[
        \boxed{
        \widehat{V}_N
        =
        e^{-r(T-t)}\overline{Y}_N.
        }
    \]

    \vspace{0.3cm}

    Son erreur standard est :

    \[
        \boxed{
        \widehat{\mathrm{SE}}
        =
        e^{-r(T-t)}
        \frac{\widehat{\sigma}_N}{\sqrt{N}}
        }
    \]


\end{frame}


%------------------------------------------------
%------------------------------------------------
\begin{frame}
    \frametitle{Intervalle de confiance à 95\,\%}

    Pour un nombre de simulations $N$ suffisamment grand,
    le théorème central limite donne approximativement :

    \[
        \frac{\widehat{V}_N-V_t}
        {\widehat{\mathrm{SE}}}
        \sim \mathcal{N}(0,1).
    \]

    \vspace{0.35cm}

    Comme une variable normale centrée réduite appartient à
    $[-1.96,1.96]$ avec une probabilité proche de $95\,\%$ :

    \[
        \boxed{
        \mathrm{IC}_{95\%}
        =
        \left[
            \widehat{V}_N-1.96\,\widehat{\mathrm{SE}},
            \;
            \widehat{V}_N+1.96\,\widehat{\mathrm{SE}}
        \right].
        }
    \]


\end{frame}


%------------------------------------------------
\begin{frame}
    \frametitle{Réduction de Variance : variables antithétiques}

    L'idée générale est de construire deux variables aléatoires
    ayant la même loi, mais négativement corrélées.

    \vspace{0.3cm}

    On note les deux payoffs :

    \[
        X^+
        \qquad\text{et}\qquad
        X^-.
    \]

    Ils vérifient :

    \[
        \mathbb{E}[X^+]
        =
        \mathbb{E}[X^-],
        \qquad
        \mathrm{Cov}(X^+,X^-)<0.
    \]

    \vspace{0.3cm}

    Au lieu d'utiliser un seul payoff, on utilise leur moyenne :

    \[
        \boxed{
        X^{anti}
        =
        \frac{X^+ + X^-}{2}.
        }
    \]

    \vspace{0.25cm}

\end{frame}






%------------------------------------------------
\begin{frame}
    \frametitle{Réduction de variance : variables antithétiques}


    Dans le modèle de Black--Scholes :

    \[
        S_T^{+}
        =
        S_t
        \exp\left[
            \left(r-\frac{\sigma^2}{2}\right)(T-t)
            +
            \sigma\sqrt{T-t}\,Z
        \right],
    \]

    \[
        S_T^{-}
        =
        S_t
        \exp\left[
            \left(r-\frac{\sigma^2}{2}\right)(T-t)
            -
            \sigma\sqrt{T-t}\,Z
        \right].
    \]

    \vspace{0.25cm}

    L'estimateur utilise ensuite la moyenne des deux payoffs :

    \[
        \boxed{
        \widehat{V}_N^{\,anti}
        =
        e^{-r(T-t)}
        \frac{1}{N}
        \sum_{i=1}^{N}
        \frac{
            \Phi(S_{T,i}^{+})
            +
            \Phi(S_{T,i}^{-})
        }{2}.
        }
    \]


\end{frame}

%------------------------------------------------
\begin{frame}
    \frametitle{Réduction de variance : variable de contrôle}

    L'idée est d'utiliser une variable aléatoire $X$ dont l'espérance
    est connue, et qui est fortement corrélée avec le payoff $Y$.

    \vspace{0.3cm}

    On suppose connaître :

    \[
        \mu_X = \mathbb{E}[X].
    \]

    \vspace{0.3cm}

    On construit alors une nouvelle variable :

    \[
        \boxed{
        Y^{ctrl}
        =
        Y
        -
        \beta
        \left(
            X-\mu_X
        \right).
        }
    \]

    \vspace{0.3cm}

    Comme $\mathbb{E}[X-\mu_X]=0$, l'espérance ne change pas :

    \[
        \mathbb{E}[Y^{ctrl}]
        =
        \mathbb{E}[Y].
    \]

\end{frame}

%------------------------------------------------
\begin{frame}
    \frametitle{Pourquoi la variance peut diminuer ?}

    La variance de la variable corrigée est :

    \[
        \mathrm{Var}(Y^{ctrl})
        =
        \mathrm{Var}(Y)
        +
        \beta^2\mathrm{Var}(X)
        -
        2\beta\,\mathrm{Cov}(Y,X).
    \]

    \vspace{0.35cm}

    Si $X$ est fortement corrélée avec $Y$, alors une partie des
    fluctuations de $Y$ peut être expliquée par les fluctuations de $X$.

    \vspace{0.35cm}

    Le coefficient optimal est :

    \[
        \boxed{
        \beta^\star
        =
        \frac{\mathrm{Cov}(Y,X)}
        {\mathrm{Var}(X)}.
        }
    \]

    \vspace{0.25cm}

    \centering
    \small
    Plus la corrélation entre $X$ et $Y$ est forte,
    plus la réduction de variance peut être importante.

\end{frame}

%------------------------------------------------
\begin{frame}
    \frametitle{Estimateur avec variable de contrôle}

    À partir de $N$ simulations, on calcule :

    \[
        Y_i^{ctrl}
        =
        Y_i
        -
        \beta
        \left(
            X_i-\mu_X
        \right).
    \]

    \vspace{0.3cm}

    L'estimateur Monte Carlo devient :

    \[
        \boxed{
        \widehat{V}_N^{ctrl}
        =
        \frac{1}{N}
        \sum_{i=1}^{N}
        Y_i^{ctrl}.
        }
    \]

    \vspace{0.3cm}

    En pratique, $\beta$ peut être estimé à partir des simulations :

    \[
        \widehat{\beta}
        =
        \frac{
            \widehat{\mathrm{Cov}}(Y,X)
        }{
            \widehat{\mathrm{Var}}(X)
        }.
    \]

\end{frame}

%------------------------------------------------
\begin{frame}
    \frametitle{Application au call de Black--Scholes}

    Pour un call européen, le payoff actualisé :

    \[
        Y_i
=
e^{-r(T-t)}
\left(S_T^{(i)}-K\right)^+
    \]

    \vspace{0.25cm}

    Une variable de contrôle naturelle est le sous-jacent actualisé :

    \[
        X_i
=
e^{-r(T-t)}S_T^{(i)}.
    \]

    \vspace{0.25cm}

    Sous la mesure risque-neutre, sans dividendes :

    \[
        \boxed{
        \mathbb{E}^{\mathbb{Q}}[X_i\mid\mathcal{F}_t]
        =
        S_t.
        }
    \]

    \vspace{0.25cm}

    On utilise donc :

    \[
        \boxed{
        Y_i^{ctrl}
        =
        Y_i
        -
        \beta
        \left(
            X_i-S_t
        \right).
        }
    \]

\end{frame}





%------------------------------------------------
\begin{frame}
    \frametitle{Pourquoi paralléliser Monte Carlo ?}

    Les simulations sont indépendantes :

    \[
        Z_i \longrightarrow S_{T,i} \longrightarrow (Y_i,X_i),
        \qquad i=1,\ldots,N.
    \]

    \vspace{0.4cm}

    Elles peuvent donc être réparties entre plusieurs cœurs du processeur.

    \vspace{0.4cm}

    \[
        \boxed{
        \text{Calculs locaux en parallèle}
        \quad\longrightarrow\quad
        \text{Fusion statistique}
        }
    \]
\end{frame}

%------------------------------------------------
\begin{frame}
    \frametitle{Répartition du travail avec OpenMP}

    OpenMP crée une équipe de threads et partage les simulations entre eux.

    \vspace{0.4cm}

    \[
    \begin{array}{c|c}
        \text{Thread} & \text{Simulations traitées} \\
        \hline
        0 & 1,\ldots,N_0 \\
        1 & N_0+1,\ldots,N_1 \\
        \vdots & \vdots \\
        p-1 & N_{p-2}+1,\ldots,N
    \end{array}
    \]

    \vspace{0.4cm}

    Chaque simulation est exécutée une seule fois.
\end{frame}

%------------------------------------------------
\begin{frame}
    \frametitle{Des données locales à chaque thread}

    Chaque thread possède :

    \begin{itemize}
        \item son propre générateur pseudo-aléatoire ;
        \item sa moyenne locale ;
        \item ses quantités locales de variance et de covariance.
    \end{itemize}

    \vspace{0.4cm}

    \[
        \mathcal{S}_k
        =
        \left(n_k,\overline{Y}_k,\overline{X}_k,
        M_{2,Y,k},M_{2,X,k},C_{YX,k}\right).
    \]

    Les threads travaillent ainsi sans modifier simultanément les mêmes données.
\end{frame}

%------------------------------------------------
\begin{frame}
    \frametitle{Fusion des statistiques de Welford}

    Pour fusionner deux groupes $A$ et $B$, on pose :

    \[
        n=n_A+n_B,
        \qquad
        \delta_Y=\overline{Y}_B-\overline{Y}_A.
    \]

    La moyenne et la dispersion globales sont :

    \[
        \boxed{\overline{Y}
        =\overline{Y}_A+\delta_Y\frac{n_B}{n}}
    \]

    \[
        \boxed{M_{2,Y}
        =M_{2,Y,A}+M_{2,Y,B}
        +\delta_Y^2\frac{n_A n_B}{n}}
    \]

    Le dernier terme mesure l'écart entre les moyennes des deux groupes.
\end{frame}

%------------------------------------------------
\begin{frame}
    \frametitle{Origine de la formule de fusion}

    Pour $i\in A$, on décompose :

    \[
        Y_i-\overline{Y}
        =(Y_i-\overline{Y}_A)+(\overline{Y}_A-\overline{Y}).
    \]

    Comme

    \[
        \sum_{i\in A}(Y_i-\overline{Y}_A)=0,
    \]

    les termes croisés disparaissent. Ainsi :

    \[
        M_{2,Y}=M_{2,Y,A}+M_{2,Y,B}
        +n_A(\overline{Y}_A-\overline{Y})^2
        +n_B(\overline{Y}_B-\overline{Y})^2.
    \]

    En remplaçant les deux écarts à la moyenne globale, on retrouve :

    \[
        \boxed{M_{2,Y}=M_{2,Y,A}+M_{2,Y,B}
        +\delta_Y^2\frac{n_A n_B}{n}}.
    \]
\end{frame}

%------------------------------------------------
\begin{frame}
    \frametitle{Fusion de la covariance}

    En posant $\delta_X=\overline{X}_B-\overline{X}_A$ :

    \[
        \boxed{
        C_{YX}
        =C_{YX,A}+C_{YX,B}
        +\delta_Y\delta_X\frac{n_A n_B}{n}
        }
    \]

    Après la fusion de tous les threads :

    \[
        \widehat{\mathrm{Var}}(Y)=\frac{M_{2,Y}}{n-1},
        \qquad
        \widehat{\mathrm{Cov}}(Y,X)=\frac{C_{YX}}{n-1}.
    \]

    Les statistiques globales sont obtenues sans stocker les simulations.
\end{frame}

%------------------------------------------------
\begin{frame}
    \frametitle{Synchronisation des threads}

    \begin{block}{\texttt{nowait}}
        Un thread ayant terminé ses simulations peut poursuivre sans attendre
        la fin des autres threads.
    \end{block}

    \begin{block}{\texttt{critical}}
        La fusion avec l'état global est exécutée par un seul thread à la fois.
        Les autres threads attendent brièvement.
    \end{block}

    \[
        \underbrace{\text{nombreuses simulations parallèles}}_{N}
        \quad\longrightarrow\quad
        \underbrace{\text{quelques fusions}}_{p}.
    \]
\end{frame}

%------------------------------------------------
\begin{frame}
    \frametitle{Bilan de la parallélisation}

    \begin{itemize}
        \item Les simulations sont réparties entre les threads.
        \item Chaque thread utilise un générateur et un état Welford privés.
        \item Les états locaux sont fusionnés mathématiquement.
        \item Une courte section critique protège la fusion globale.
    \end{itemize}

    \vspace{0.4cm}

    L'accélération apportée par $p$ threads se mesure par :

    \[
        \boxed{\mathrm{Speedup}(p)=\frac{T_1}{T_p}}
    \]

    Elle reste généralement inférieure à $p$ en raison des synchronisations
    et des limites matérielles.
\end{frame}

\end{document}
